\chapter{Haskell 初体验}\label{haskell-ux521dux4f53ux9a8c}

本章先用一个最小脚本熟悉 Haskell 的基本语法与 GHCi（1.1 节）；再用同一个求值器演示三种
写法：Python 的 \passthrough{\lstinline!isinstance!} 派发、抽象基类、Haskell 的 \passthrough{\lstinline!data!} 加模式匹配（1.2 节），并在
1.3 节比较三种实现的取舍，引出“表达问题”。1.4 节把比较延伸到状态处理、C/Go/Scala 与各家 REPL；1.5 节看函数式语言内部的差异与 Haskell 的近亲；1.6 节讲惰性求值与副作用。读完本章应能区分：Haskell 与 Python 的差别主要不在语法，而在类型系统的检查能力与默认
求值方式。语法细节、类型类规则与范畴论留到后续章节。

\section{脚本与 GHCi}\label{ux811aux672cux4e0e-ghci}

\subsection{脚本}\label{ux811aux672c}

Haskell 程序由一组\textbf{定义}组成：先写类型签名，再写实现。一个最小的脚本：

\begin{lstlisting}[language=Haskell]
-- Comments start with two dashes.
double :: Int -> Int
double x = x * 2

-- The entry point: an IO () action.
main :: IO ()
main = print (double 21)
\end{lstlisting}

\passthrough{\lstinline!double :: Int -> Int!} 是签名，读作“double 把 \passthrough{\lstinline!Int!} 映射到 \passthrough{\lstinline!Int!}”；\passthrough{\lstinline!main :: IO ()!} 是程序
入口，此处只需知道它是“执行后返回 \passthrough{\lstinline!()!} 的动作”，\passthrough{\lstinline!IO!} 的类型与组合规则在工程各章随用随讲。运行：

\begin{lstlisting}
$ runghc code/src/Ch01/Hello.hs
42
\end{lstlisting}

这个例子让初学者快速了解 Haskell 的基本语法。本书不追求语法细节的完整覆盖。缩进规则、\passthrough{\lstinline!let!} 的完整规则、模块头语法等，第 2--4 章按“够用”原则引入；写代码时以编译器报错为准，具体语法需要时查手册即可。

\subsection{GHCi 用法与注意事项}\label{ghci-ux7528ux6cd5ux4e0eux6ce8ux610fux4e8bux9879}

交互地试一行定义、查一个类型，用 \textbf{GHCi}（Haskell 的 REPL）：

\begin{lstlisting}
ghci> :l code/src/Ch01/Hello.hs
ghci> double 21
42
ghci> :t double
double :: Int -> Int
\end{lstlisting}

在 GHCi 中，每行以 \passthrough{\lstinline!ghci>!} 开头等待输入；以 \passthrough{\lstinline!:!} 开头的词是命令而非表达式：\passthrough{\lstinline!:l!}（load）加载
文件，\passthrough{\lstinline!:t!}（type）查类型，\passthrough{\lstinline!:i!}（info）看定义，\passthrough{\lstinline!:r!}（reload）重新加载。上面第三行输出
\passthrough{\lstinline!double :: Int -> Int!}：\passthrough{\lstinline!::!} 读作“具有类型”，这一行是签名，不是运行结果。GHCi 对每个输入
表达式求值后立即打印其值，因此 \passthrough{\lstinline!double 21!} 的下一行就是 \passthrough{\lstinline!42!}。

使用中有几件事需要注意。GHCi 里的 \passthrough{\lstinline!let!} 是定义不是赋值：重复 \passthrough{\lstinline!let x = 1!} 再 \passthrough{\lstinline!let x = 2!} 是
遮蔽（shadowing），两个 \passthrough{\lstinline!x!} 并存、旧定义不受影响，1.4.1 节“没有可变变量，只有定义”在 REPL
里同样成立。跨行的定义要包在 \passthrough{\lstinline!:\{ ... :\}!} 里（或 \passthrough{\lstinline!:set +m!}），否则函数体写不下一行时直接
回车，GHCi 会提前对已输入的部分求值。函数本身不能打印：\passthrough{\lstinline!ghci> id!} 报错
\passthrough{\lstinline!No instance for (Show (a0 -> a0))!}，因为函数没有显示方式。惰性只推迟计算、不推迟打印：
\passthrough{\lstinline![1..]!} 不会挂起，但 \passthrough{\lstinline!print [1..]!} 会无限输出（1.6 节）。

\section{一个求值器}\label{ux4e00ux4e2aux6c42ux503cux5668}

任务是：给定一棵算术表达式树，求出它的值。

用 Python 实现，最直接的写法是四个只保存字段的类，加一个用 \passthrough{\lstinline!isinstance!} 判断形状的函数。

\subsection{第一种实现：isinstance 派发}\label{ux7b2cux4e00ux79cdux5b9eux73b0isinstance-ux6d3eux53d1}

\begin{lstlisting}[language=Python]
class Lit:
    def __init__(self, n): self.n = n

class Add:
    def __init__(self, a, b): self.a, self.b = a, b

class Mul:
    def __init__(self, a, b): self.a, self.b = a, b

class Neg:
    def __init__(self, a): self.a = a

def evaluate(e):
    if isinstance(e, Lit): return e.n
    if isinstance(e, Add): return evaluate(e.a) + evaluate(e.b)
    if isinstance(e, Mul): return evaluate(e.a) * evaluate(e.b)
    if isinstance(e, Neg): return -evaluate(e.a)
    raise TypeError(f"unknown expression: {e!r}")

print(evaluate(Mul(Lit(6), Add(Lit(2), Lit(3)))))   # 30
\end{lstlisting}

输出 \passthrough{\lstinline!30!}（\passthrough{\lstinline!code/python/eval\_naive.py!}，CI 中执行）。

这段代码的结构与后面的 Haskell 版本基本相同：都是一个函数按形状分支。区别不在结构，而在谁
保证分支完整。\passthrough{\lstinline!isinstance!} 链依赖运行时反射，Python 并不知道 \passthrough{\lstinline!Expr!} 有几种形状（这里也没有
\passthrough{\lstinline!Expr!} 这个类型），因此最后一行 \passthrough{\lstinline!raise TypeError!} 承担了本应由编译器承担的工作。

\subsection{第二种实现：抽象基类}\label{ux7b2cux4e8cux79cdux5b9eux73b0ux62bdux8c61ux57faux7c7b}

Python 也可以写成抽象基类加方法分派，\passthrough{\lstinline!if!} 链随之消失：

\begin{lstlisting}[language=Python]
from abc import ABC, abstractmethod


class Node(ABC):
    @abstractmethod
    def eval(self) -> int: ...

class Lit(Node):
    def __init__(self, n: int) -> None: self.n = n

    def eval(self) -> int:
        return self.n

class Add(Node):
    def __init__(self, a: Node, b: Node) -> None: self.a, self.b = a, b

    def eval(self) -> int:
        return self.a.eval() + self.b.eval()
\end{lstlisting}

\passthrough{\lstinline!Mul!}、\passthrough{\lstinline!Neg!} 同理，完整文件在 \passthrough{\lstinline!code/python/eval\_oop.py!}。

\passthrough{\lstinline!Mul(Lit(6), Add(Lit(2), Lit(3))).eval()!} 同样得到 \passthrough{\lstinline!30!}。每个类只处理自己那一种形状，增加
一种节点只需写一个类，其余代码不变。这是相对第一种实现的实质改进。

但这种实现仍有代价，只是代价转移到别处：

\begin{itemize}
\tightlist
\item
  \textbf{一个操作被拆分到四个类中。} \passthrough{\lstinline!eval!} 的定义分散在各处，阅读时需要来回跳转；修改语义
  （例如除法取整还是取浮点）需逐一修改。
\item
  \textbf{增加一种操作仍需改动每个类。} 1.3.3 节增加 \passthrough{\lstinline!pretty!} 时，四个类各要写一个方法。
\item
  \textbf{“共有哪些形状”仍无一处完整写出。} \passthrough{\lstinline!Node!} 只知道自己是基类，不知道有多少子类；
  \passthrough{\lstinline!isinstance!} 与 \passthrough{\lstinline!issubclass!} 也提供不了这一信息。
\item
  \textbf{类型注解无法补足。} mypy 检查的是“每个类都实现了 \passthrough{\lstinline!eval!}”，不是“处理了所有形状”，
  其能力边界与抽象基类相同。
\end{itemize}

第二种实现得到的是“增加节点容易”与“以对象组织代码”，没有获得“编译器掌握全部情况”。

\subsection{Haskell 实现}\label{haskell-ux5b9eux73b0}

先定义表达式有哪些形状：

\begin{lstlisting}[language=Haskell]
data Expr
    = Lit Int
    | Add Expr Expr
    | Mul Expr Expr
    | Neg Expr
    deriving (Eq, Show)
\end{lstlisting}

一行 \passthrough{\lstinline!data!} 声明了四种情况。\passthrough{\lstinline!Lit Int!} 是叶子，\passthrough{\lstinline!Add Expr Expr!} 是两个子表达式的组合，
\passthrough{\lstinline!deriving (Eq, Show)!} 由编译器生成相等判断与打印函数。Python 中若要 \passthrough{\lstinline!\_\_eq\_\_!} 与 \passthrough{\lstinline!\_\_repr\_\_!}
通常需自行编写（\passthrough{\lstinline!@dataclass!} 可代为生成，但那是标准库的功能，不属于语言本身）。

求值：

\begin{lstlisting}[language=Haskell]
eval :: Expr -> Int
eval (Lit n)   = n
eval (Add a b) = eval a + eval b
eval (Mul a b) = eval a * eval b
eval (Neg a)   = negate (eval a)
\end{lstlisting}

四个方程，每行左侧一个形状，右侧一个结果。没有 \passthrough{\lstinline!return!}，没有 \passthrough{\lstinline!isinstance!}，也没有兜底异常。
使用方式：

\begin{lstlisting}[language=Haskell]
demoExpr :: Expr
demoExpr = Mul (Lit 6) (Add (Lit 2) (Lit 3))
\end{lstlisting}

\passthrough{\lstinline!eval demoExpr!} 得到 \passthrough{\lstinline!30!}。这些代码位于 \passthrough{\lstinline!code/!}，每次提交由 GHC 编译、由 hspec 断言，书中
引用的每一行都可运行。

三种写法并排：

{\small\setlength{\tabcolsep}{4pt}\def\LTcaptype{table} % do not increment counter
\begin{longtable}[]{@{}
  l
  l
  l
  l@{}}
\caption{三种写法的对照}\label{tab:1-1}\\
\toprule\noalign{}
\textbf{} & \textbf{形状写在哪} & \textbf{一个操作写在几个地方} & \textbf{漏了一种，谁知道} \\
\midrule\noalign{}
\endhead
\bottomrule\noalign{}
\endlastfoot
Python 第一种 & 四个哑类 + \passthrough{\lstinline!isinstance!} 链 & 一处，靠运行时反射 & 运行时抛 \passthrough{\lstinline!TypeError!} \\
Python 第二种 & 四个类，各自一个方法 & 四处，每个类一处 & 实例化时，抽象基类拦截 \\
Haskell & 一处 \passthrough{\lstinline!data!} 声明 & 一处，一个函数 & 编译时，并给出名字 \\
\end{longtable}
}

表 1.1 有两点需要注意。第一，Python 两种写法的差别只在第二列，即操作写在一处还是四处，这属于代码组织问题。第二，Haskell 多出的能力在第三列，即穷尽性由编译器保证，这属于检查问题。
两者正交：操作放在哪里，以及谁负责检查完整性。

至于哪种更好，取决于代码将来的修改方向。1.3.2 与 1.3.3 节分别沿两个方向各走一遍，
1.3.4 节讨论这一取舍。

\section{三种实现的比较}\label{ux4e09ux79cdux5b9eux73b0ux7684ux6bd4ux8f83}

\subsection{函数体是表达式}\label{ux51fdux6570ux4f53ux662fux8868ux8fbeux5f0f}

Python 的函数体是一串语句：赋值、分支、循环、返回。Haskell 的函数体是一个表达式，
\passthrough{\lstinline!if!}、\passthrough{\lstinline!case!}、\passthrough{\lstinline!let!} 都是表达式，都有值。

因此 Haskell 没有 \passthrough{\lstinline!return!} 语句，函数的值就是最后一个表达式的值。语言也不规定
\passthrough{\lstinline!eval (Add a b)!} 中先算 \passthrough{\lstinline!a!} 还是先算 \passthrough{\lstinline!b!}。

这一规定的实际后果是：函数内部无法修改状态。需要修改时，必须把新状态作为返回值传递。
第 9 章的 State 单子把这一过程封装起来，此处暂不展开。

\subsection{增加节点}\label{ux589eux52a0ux8282ux70b9}

假设需要支持除法。以第二种实现为对照（它强于第一种），只需增加一个类：

\begin{lstlisting}[language=Python]
class Div(Node):
    def __init__(self, a: Node, b: Node) -> None:
        self.a, self.b = a, b

    def eval(self) -> int:
        return self.a.eval() // self.b.eval()
\end{lstlisting}

其余代码不变，\passthrough{\lstinline!Div(Lit(7), Lit(2)).eval()!} 直接得到 \passthrough{\lstinline!3!}。新类自动成为 \passthrough{\lstinline!Node!} 的子类，
也自动被 \passthrough{\lstinline!expr.eval()!} 使用。

Haskell 中则需要修改数据类型本身：

\begin{lstlisting}[language=Haskell]
data Expr = Lit Int | Add Expr Expr | Mul Expr Expr | Neg Expr | Div Expr Expr
\end{lstlisting}

然后用 \passthrough{\lstinline!-Wall!} 编译，GHC 会指出哪些函数遗漏了这一分支。\passthrough{\lstinline!-Wall!} 是 GHC 的编译选项：\passthrough{\lstinline!-W!} 是警告
选项的前缀，\passthrough{\lstinline!-Wall!} 打开其中常用的一批（并非全部，例如 \passthrough{\lstinline!-Wmissing-export-lists!} 不在其内），
负责这一项的是 \passthrough{\lstinline!-Wincomplete-patterns!}。开启方式有三种：命令行 \passthrough{\lstinline!ghc -Wall!}、工程文件里的
\passthrough{\lstinline!ghc-options: -Wall!}、源文件顶部的 \passthrough{\lstinline!\{-\# OPTIONS\_GHC -Wall \#-\}!}。以下是本机实际输出：

\begin{lstlisting}
ExpressionProblem.hs:12:1: warning: [GHC-62161] [-Wincomplete-patterns]
        Patterns of type 'Expr' not matched: Div _ _
12 | eval (Lit n)   = n

ExpressionProblem.hs:18:1: warning: [GHC-62161] [-Wincomplete-patterns]
        Patterns of type 'Expr' not matched: Div _ _
18 | pretty (Lit n)   = show n
\end{lstlisting}

两条警告分别对应 \passthrough{\lstinline!eval!} 与 \passthrough{\lstinline!pretty!}，既指出遗漏的存在，也指出遗漏所在的函数与构造子。
这是“类型封闭”的结果：\passthrough{\lstinline!data!} 声明处列出全部构造子，编译器得以逐个核对。

在增加节点这一点上，Python 的实现更省力。下一节改变方向。

\subsection{增加操作}\label{ux589eux52a0ux64cdux4f5c}

不改动节点，增加一个操作：把表达式打印成中缀字符串（\passthrough{\lstinline!pretty!}）。它不改变求值结果，只是
以另一种形式表示同一棵树。

Haskell 中只需增加一个函数：

\begin{lstlisting}[language=Haskell]
pretty :: Expr -> String
pretty (Lit n)   = show n
pretty (Add a b) = "(" ++ pretty a ++ " + " ++ pretty b ++ ")"
pretty (Mul a b) = "(" ++ pretty a ++ " * " ++ pretty b ++ ")"
pretty (Neg a)   = "-" ++ pretty a
\end{lstlisting}

\passthrough{\lstinline!pretty demoExpr!} 得到 \passthrough{\lstinline!"(6 * (2 + 3))"!}。\passthrough{\lstinline!Expr!} 与 \passthrough{\lstinline!eval!} 均未改动。

Python 中需要在 \passthrough{\lstinline!Node!} 上增加抽象方法，再回到四个类逐一实现：

\begin{lstlisting}[language=Python]
class Node(ABC):
    @abstractmethod
    def eval(self) -> int: ...
    @abstractmethod
    def pretty(self) -> str: ...


class Lit(Node):
    def __init__(self, n: int) -> None:
        self.n = n

    def eval(self) -> int:
        return self.n

    def pretty(self) -> str:
        return str(self.n)
# Add / Mul / Neg 各来一遍
\end{lstlisting}

遗漏的实现会被抽象基类拦截，但拦截发生在实例化时：

\begin{lstlisting}
TypeError: Can't instantiate abstract class Broken without an implementation for abstract method 'eval'
\end{lstlisting}

（该输出由 \passthrough{\lstinline!code/python/eval\_oop.py!} 实际产生。）它的能力与边界都需注意：它保证“每个类都
实现了 \passthrough{\lstinline!eval!}“，不保证”处理了所有类“；且错误在 \passthrough{\lstinline!Broken()!} 这一行触发，若该类只在很少的
路径上被构造，问题可能延迟到运行期才暴露。

在增加操作这一点上，Haskell 的实现更省力。需要补充的是，第一种实现在这一方向上的改动量也不大：
再写一个 \passthrough{\lstinline!pretty(e)!} 函数即可。但它仍需自行编写 \passthrough{\lstinline!isinstance!} 链并兜底，且与 \passthrough{\lstinline!evaluate!} 一样
没有任何静态保证。

\subsection{表达问题}\label{ux8868ux8fbeux95eeux9898}

两个方向的改动量并列如下：

{\small\setlength{\tabcolsep}{4pt}\def\LTcaptype{table} % do not increment counter
\begin{longtable}[]{@{}
  l
  l
  l@{}}
\caption{两类扩展的代价}\label{tab:1-2}\\
\toprule\noalign{}
\textbf{} & \textbf{增加一种节点} & \textbf{增加一种操作} \\
\midrule\noalign{}
\endhead
\bottomrule\noalign{}
\endlastfoot
Python 的类 & 增加一个类，其余不改 & 每个类都要改 \\
Haskell 的 ADT & 修改定义 + 所有匹配它的函数（编译器列出清单） & 增加一个函数 \\
\end{longtable}
}

表 1.2 列出的不是优劣，而是选择：代码将来主要沿哪个方向扩展。

\begin{itemize}
\tightlist
\item
  节点固定、操作会不断增加（求值、打印、类型检查、优化、编译为字节码），适合用 ADT。
  实际编译器多数如此，GHC 本身也是。
\item
  操作固定、类型会不断增加（例如一个“保存到磁盘”的接口，将来支持更多文件格式），适合用对象。
\end{itemize}

这一问题称为\textbf{表达问题}（the expression problem），没有无代价的解法，只能决定把代价放在
哪一侧。Haskell 与 Python 恰好各在一侧。（名称来历见本章末参考文献。）

Haskell 也能写成对象的形式。把“每个节点自带方法”直接迁移过来，最简单的形式是：

\begin{lstlisting}[language=Haskell]
data Expr = Expr { eval :: Int }
\end{lstlisting}

节点在此成为记录：\passthrough{\lstinline!eval!} 是构造时算好的一个 \passthrough{\lstinline!Int!}，形状信息在构造时即被吸收。这条路
（final style 或 Church 编码）可用，也有其适用场合，但它使“同一棵树、多种解释”不再成立：
值的解释在构造时已经确定。

\subsection{函数作为值}\label{ux51fdux6570ux4f5cux4e3aux503c}

Python 同样支持\textbf{高阶函数}（higher-order function）：\passthrough{\lstinline!map!}、\passthrough{\lstinline!filter!}、\passthrough{\lstinline!sorted(key=...)!}
都是常用形式。Haskell 的差别不在“是否支持”，而在这是默认的组织方式。

例：把一串数各自平方。

\begin{lstlisting}[language=Haskell]
squaresOf :: [Int] -> [Int]
squaresOf = map (^ 2)
\end{lstlisting}

\passthrough{\lstinline!map (\^ 2)!} 本身即函数 \passthrough{\lstinline![Int] -> [Int]!}，右侧无须写参数。Python 中对应写法是列表推导
\passthrough{\lstinline![x*x for x in xs]!}，或 \passthrough{\lstinline!lambda!} 加 \passthrough{\lstinline!list(map(...))!}；后者可写但不常见，因为嵌套较深。
语言的默认表达方式会影响代码的组织形式。

\section{与命令式语言比较}\label{ux4e0eux547dux4ee4ux5f0fux8bedux8a00ux6bd4ux8f83}

\subsection{状态}\label{ux72b6ux6001}

把列表求和，Python 的常见写法：

\begin{lstlisting}[language=Python]
total = 0
for x in xs:
    total += x
\end{lstlisting}

包含三个步骤：初始化、循环、累加。\passthrough{\lstinline!total!} 在迭代中多次改变取值，其正确性依赖两个前提：
循环体按顺序执行，且每轮都更新 \passthrough{\lstinline!total!}。

Haskell 写作 \passthrough{\lstinline!sum xs!}，不使用循环。数学上，求和是把 \passthrough{\lstinline!+!} 插入列表元素之间：

\begin{lstlisting}[language=Haskell]
foldr (+) 0 [1,2,3]   -- 1 + (2 + (3 + 0))
\end{lstlisting}

（\passthrough{\lstinline!base!} 为性能考虑实际使用严格左折叠 \passthrough{\lstinline!foldl' (+) 0!}，结果相同；必须如此的原因见第 2 章
关于惰性的讨论。）

这里没有计数器，没有“当前值”，也没有“第几次迭代”。

需要区分的是：\textbf{Haskell 没有可变变量，只有定义}。\passthrough{\lstinline!x = 5!} 不是把 5 存入 \passthrough{\lstinline!x!}，而是声明从此
\passthrough{\lstinline!x!} 表示 5。\passthrough{\lstinline!x = x + 1!} 能通过类型检查，但它不是赋值，而是递归定义：这个等式没有有限解，
求值得到 \(\bot\)。需要可变的量时，有两种做法：作为参数显式传递，或用单子封装传递过程（第 9
章）。前者属于数学，后者属于工程，本书两者都会涉及。

C 程序员会问：用指针、\passthrough{\lstinline!for!} 循环与可变数组不也可以吗？可以，但代价是：无法从函数签名判断它
是否会修改传入的参数。Haskell 的类型签名提供了这一保证：\passthrough{\lstinline!eval :: Expr -> Int!} 除了算出那个
\passthrough{\lstinline!Int!}，不做别的事。

\subsection{与 C、Go、Scala 比较}\label{ux4e0e-cgoscala-ux6bd4ux8f83}

1.3 节比较了 Python 与 Haskell。把 C、Go、Scala 一并纳入，问题仍是同一组：形状如何表达，
谁能增加新形状，遗漏分支时由谁在何时报告。各语言的答案见表 1.3。

{\footnotesize\setlength{\tabcolsep}{4pt}\def\LTcaptype{table} % do not increment counter
\begin{longtable}[]{@{}
  l
  l
  l
  l@{}}
\caption{C、Go、Scala、Haskell 的对照}\label{tab:1-3}\\
\toprule\noalign{}
\textbf{语言} & \textbf{形状怎么表达} & \textbf{谁能加新形状} & \textbf{漏写分支会怎样} \\
\midrule\noalign{}
\endhead
\bottomrule\noalign{}
\endlastfoot
C & 手写的 \passthrough{\lstinline!enum!} / \passthrough{\lstinline!union!} / \passthrough{\lstinline!switch!} & 任何人（改头文件） & 编译通过；\passthrough{\lstinline!-Wswitch!} 只是提醒 \\
Go & 接口 + type switch & 任何人（实现方法即可） & 编译通过；没有穷尽性检查 \\
Scala & \passthrough{\lstinline!sealed trait!} + \passthrough{\lstinline!case class!} & 只有同文件内（\passthrough{\lstinline!sealed!}） & 编译警告 \\
Haskell & \passthrough{\lstinline!data Expr = Lit Int \\| …!} & 只有定义处（封闭） & 编译警告，并列出遗漏的名字 \\
\end{longtable}
}

\textbf{C} 使用手写的 ADT：\passthrough{\lstinline!enum!} 是构造子表，\passthrough{\lstinline!union!} 是负载，\passthrough{\lstinline!switch!} 是\textbf{模式匹配}（pattern matching）。技术上它是
封闭的，但任何代码都可以修改头文件增加一个 tag，其他编译单元不会因此重新检查。遗漏分支时
编译器不报告（\passthrough{\lstinline!-Wswitch!} 仅在开启时提醒 enum 少列）。这与其设计定位一致：编译器不替程序员
决定内存布局。

\textbf{Go} 的 \passthrough{\lstinline!interface!} 是开放集合：任何类型只要实现相应方法即自动满足，包括尚未编写的类型。
成员在原则上不可枚举，\passthrough{\lstinline!type switch!} 因此没有穷尽性可言：遗漏分支走 \passthrough{\lstinline!default!}，或者什么
也不做。Go 用可检查性换取了可扩展性，这一取舍符合其设计目标。

\textbf{Scala} 的 \passthrough{\lstinline!sealed trait!} 将子类限制在同一文件内，非穷尽 \passthrough{\lstinline!match!} 会给出警告，与 GHC 属于
同一类。它的差异在其他方面：默认严格求值、副作用不受限制、类型系统更复杂。下一节讨论。

\textbf{Haskell} 的情况已在 1.3.2 节说明：封闭，且由编译器检查。

需要指出的是，上表实际混合了三个独立的问题：

\begin{enumerate}
\def\labelenumi{\arabic{enumi}.}
\tightlist
\item
  \textbf{能否穷举？} 形状表能否在一处列完。C 能，但依赖纪律；Go 原则上不能；Scala 与 Haskell
  能，且由语言强制。
\item
  \textbf{由谁在何时检查？} 编译期（Scala、Haskell）、运行时（Python 的抽象基类在实例化时），
  或不做检查（C、Go）。
\item
  \textbf{谁能扩展这张表？} 任何位置（C 的头文件、Go 的接口、Python 的子类），或仅限一处
  （Scala 的 \passthrough{\lstinline!sealed!}、Haskell 的 \passthrough{\lstinline!data!}）。
\end{enumerate}

三者可以独立组合。因此“Haskell 的类型系统更强”这一说法信息量有限；更准确的表述是：Haskell
选择了“可穷举 + 编译期检查 + 单点扩展”，并接受了相应的代价。

最后是 \passthrough{\lstinline!null!}，它不属于上述三项中的任何一项，而是第四件事：类型是否把可能性写全。C 的
\passthrough{\lstinline!NULL!}、Go 的 \passthrough{\lstinline!nil!}、Java 的 \passthrough{\lstinline!null!} 使 \passthrough{\lstinline!String!} 同时表示“有一个字符串”与“什么都没有”，
类型未写全，因此每个调用点都要补一个判断。Haskell 把这一情形写入类型：可能缺失即 \passthrough{\lstinline!Maybe a!}，
取值为 \passthrough{\lstinline!Just x!} 或 \passthrough{\lstinline!Nothing!}，遗漏 \passthrough{\lstinline!Nothing!} 时编译器会提示。第 9、10 章将把 \passthrough{\lstinline!Maybe!} 作为
普通实例处理。

\subsection{各种语言的 REPL 比较}\label{ux5404ux79cdux8bedux8a00ux7684-repl-ux6bd4ux8f83}

1.1 节介绍了 GHCi 的基本用法。把不同语言的 REPL 并排看，差别同样反映语言核心：

\begin{itemize}
\tightlist
\item
  \textbf{Python} 的 REPL 同样打印表达式值，差别有两点：Python 允许 \passthrough{\lstinline!x = x + 1!} 这类修改状态的
  语句，GHCi 没有可变变量；Python 打印任何对象（函数显示为 \passthrough{\lstinline!<function id at 0x...>!} 的
  地址形式），Haskell 对没有 \passthrough{\lstinline!Show!} 实例的值直接报错。
\item
  \textbf{Node / Lisp} 的 REPL 以表达式为核心，但可以随时就地产生副作用（修改对象、打印）；
  GHCi 中副作用必须显式进入 \passthrough{\lstinline!IO!}：\passthrough{\lstinline!putStrLn "hi"!} 会执行，但对象无法就地修改。
\item
  \textbf{ML 系（SML / OCaml）} 的 REPL 与 GHCi 最接近：表达式驱动、静态类型，差别主要在求值
  策略（严格）与模块系统的组织方式。
\end{itemize}

这些差别不是交互环境的设计选择，而是语言核心（纯函数、无赋值）在 REPL 上的体现。
本书对 REPL 的介绍到此为止：日常使用以脚本运行为主，工程各章的代码都以可编译文件的形式
给出。

\section{函数式语言及其近亲}\label{ux51fdux6570ux5f0fux8bedux8a00ux53caux5176ux8fd1ux4eb2}

这一节把视野从 Haskell 扩展到整个函数式语言家族：先看家族内部三条路线的差异，再看
Haskell 的近亲。各家 REPL 的比较已放在 1.4.3 节。

\subsection{三条路线}\label{ux4e09ux6761ux8defux7ebf}

首先需要澄清一点：\textbf{函数式不等于静态类型}。Lisp 是函数式语言的源头，默认没有静态类型；
Python 支持类型注解但不强制；Haskell 的 ADT 也不是函数式的必然产物，它来自 Hindley--Milner
类型系统与 ML 家族，Scala 从中借用。函数式与静态类型是两条独立的轴。

三条路线的差别如下。

\textbf{Lisp（Scheme / Common Lisp）} 的特点是语言可塑性强。代码即列表，列表即数据，宏可以把
语法改造成任意形式。代价是抽象依赖约定与文档；求值为严格求值（call-by-value），需要惰性时
须使用 \passthrough{\lstinline!delay!} / \passthrough{\lstinline!force!} 或手写 thunk。Lisp 的方向是“语言本身可被改写”，Haskell 不在这一
方向上：它有 Template Haskell，但主流写法不以宏做抽象，而是用类型与类型类。

\textbf{Scala} 的特点是渐进落地。它运行在 JVM 上，可与 Java 库互操作，允许从命令式写法开始逐步
转向函数式。默认严格求值，需要惰性时写 \passthrough{\lstinline!lazy val!} 或按名参数（\passthrough{\lstinline!=> T!}）；Haskell 相反，默认
惰性，需要严格时用 \passthrough{\lstinline!seq!} 或 \passthrough{\lstinline"!"}。Scala 的类型系统比 Haskell 更复杂（子类型、\passthrough{\lstinline!implicit!} /
\passthrough{\lstinline!given!}、match types），代价是错误信息往往很长。其 \passthrough{\lstinline!sealed trait!} 加 \passthrough{\lstinline!case class!} 就是 ADT，
\passthrough{\lstinline!for!} 推导式即 Haskell 的 \passthrough{\lstinline!do!}，两项均借鉴自 Haskell 与 ML。

\textbf{Haskell} 同时固定了三件事：默认惰性、默认纯、以类型类实现 ad-hoc 多态。三者相互依赖：

\begin{itemize}
\tightlist
\item
  若允许随处修改状态，惰性求值的时机将影响结果，按需求值不再安全；
\item
  若默认严格，纯函数仍需规定求值顺序，抽象退化为“先算完再传”；
\item
  若没有类型类，\passthrough{\lstinline!+!} 只能绑定单一类型，或退回 C++ 式的运算符重载。
\end{itemize}

因此其他语言中的“Haskell 特性”，多数是这三条的推论，而不是彼此独立的设计选择。

\subsection{Haskell 的近亲}\label{haskell-ux7684ux8fd1ux4eb2}

内部路线之外，还有另一类关系：哪些语言直接模仿 Haskell，哪些与 Haskell 共享核心特性。
了解它们的目的不是换用它们，而是分清 Haskell 的哪些部分属于家族共享，哪些是它独有。

\textbf{PureScript} 语法与类型系统取自 Haskell（类型类、ADT、\passthrough{\lstinline!do!} 记法），语义独立设计：严格
求值，编译到 JavaScript，无运行时依赖。它在浏览器端保留了 Haskell 的抽象能力，放弃了默认
惰性，换来可预测的性能与较小的产物。

\textbf{Elm} 同样是面向前端的纯函数语言，但刻意做小：没有类型类，整个程序由固定的 Elm 架构
（model--update--view）组织。Haskell 用户会熟悉它的语法，却不会把它当作 Haskell 在浏览器上
的移植。

\textbf{Frege} 是 JVM 上的 Haskell：把 Haskell 语义搬进 Java 平台并与之互操作。项目已停止开发，
但代表“Haskell 语义、更换运行平台”这条路线；同一条路线上还有 GHCJS 与 Haste，它们把
Haskell（或子集）直接编译为 JavaScript，保留语义，代价是生成代码体积大。

\textbf{Clean} 与 Haskell 同源：两者都出自惰性纯函数语言的传统，共享惰性求值、纯函数、代数数据
类型这些核心特性，独立发展至今。Clean 用唯一性类型（uniqueness types）把副作用写进类型，
Haskell 的答案是单子（第 9 章）。

\textbf{Idris} 语法接近 Haskell，类型系统强一档：类型可以依赖值（依赖类型），能在类型中写出
“长度为 \(n\) 的列表”这类精确性质。第 3 章的 HoTT 讨论是同一方向的数学基础。对 Haskell 读者，
Idris 展示的是“把更多约束交给编译器”的极限。

对照之下，Haskell 独有的部分是默认惰性、以 \passthrough{\lstinline!IO!} 组织副作用，以及 GHC 的工程生态；语法、
ADT、类型类与纯函数则是整个家族共享的面。

\section{惰性求值与副作用}\label{ux60f0ux6027ux6c42ux503cux4e0eux526fux4f5cux7528}

\subsection{惰性}\label{ux60f0ux6027}

下面这个定义不产生计算，也不会陷入死循环：

\begin{lstlisting}[language=Haskell]
firstFive :: [Int]
firstFive = take 5 squares
  where
    squares = map (^ 2) [0 ..]   -- an infinite list; defining it costs nothing
\end{lstlisting}

\passthrough{\lstinline![0 ..]!} 是全体自然数，\passthrough{\lstinline!map (\^ 2)!} 作用其上得到另一个无限长的列表。实际被计算的只有
\passthrough{\lstinline!take 5!} 所需的前五项 \passthrough{\lstinline![0,1,4,9,16]!}，其余元素从未产生。

Python 中表达同一件事需使用生成器：\passthrough{\lstinline!(x*x for x in itertools.count())!}。区别在于 Python 的
\textbf{惰性求值}（lazy evaluation）是局部的（生成器是特殊对象），Haskell 的惰性求值是全局的
（所有值默认按需求值）。这一差别在后文会反复出现：\passthrough{\lstinline!let xs = [1..]!} 为何成立，\passthrough{\lstinline!undefined!}
为何能放进数据结构而不立即出错，以及某些线性算法为何占用大量内存。

惰性有代价。\passthrough{\lstinline!foldl (+) 0 [1..10\^7]!} 会先积累一千万层未求值的加法表达式（thunk），再一次性
求值，内存占用可达数 GB；正确写法是严格的 \passthrough{\lstinline!foldl'!}。第 2 章将详述这一问题。此处只需记住：
惰性把求值时机交给运行时，代价由使用者承担。

\subsection{副作用}\label{ux526fux4f5cux7528}

初学者常见的疑问是：既然函数不能有\textbf{副作用}（side effect），\passthrough{\lstinline!print!} 如何解释？

\passthrough{\lstinline!eval!} 的类型是 \passthrough{\lstinline!Expr -> Int!}。签名中没有 \passthrough{\lstinline!IO!}，因此 \passthrough{\lstinline!eval!} 不能打印、不能读文件、不能发起
网络请求，只能算出一个 \passthrough{\lstinline!Int!}。这不是编码约定，而是类型系统强制的结果：要产生副作用，签名中
必须出现 \passthrough{\lstinline!IO!}。

程序的入口是 \passthrough{\lstinline!main :: IO ()!}；其余 \passthrough{\lstinline!IO!} 动作在被接到它上面之前不会执行。\passthrough{\lstinline!IO!} 是类型构造子，
\passthrough{\lstinline!IO ()!} 表示“一个执行动作并返回 \passthrough{\lstinline!()!} 的值”。\passthrough{\lstinline!putStrLn "hi"!} 的类型是 \passthrough{\lstinline!String -> IO ()!}：给定字符串，返回一个动作。
动作本身是值，可以传递、组合、放入列表，但只有接到 \passthrough{\lstinline!main!} 上才会执行。

Haskell 处理副作用的方式不是禁止，而是\textbf{把副作用写进类型}，使其在签名上可见。编译器因此
可以判断哪些函数是纯的；纯函数可以被重排、缓存、并行与内联。第 8 章之后将说明，正是这一性质
使“以组合替代循环”在工程上可行。

\section{总结}\label{ux603bux7ed3}

本章做了四件事：用最小脚本与 GHCi 建立对 Haskell 的第一印象（1.1 节）；用同一个求值器
演示三种写法，引出表达问题（1.2、1.3 节）；把差异放到语言比较中看------与 C、Go、Scala 的
命令式对照、函数式家族内部的路线与近亲、各家 REPL 的差别（1.4、1.5 节）；最后讲惰性求值
与副作用（1.6 节）。

贯穿全章的有三条结论：

\begin{enumerate}
\def\labelenumi{\arabic{enumi}.}
\tightlist
\item
  \textbf{函数体是表达式}，因此没有语句序列，也没有 \passthrough{\lstinline!return!}；
\item
  \textbf{类型可以封闭}，因此遗漏一种情况会在编译期被发现，代价是增加节点时改动更多；
\item
  \textbf{副作用写在类型里}，一个函数是否纯净可以直接从签名读出。
\end{enumerate}

第 2 章起逐条展开这三点。本章遗留的问题：语法细节（第 2--4 章）、类型类的完整规则
（第 4 章）、范畴论（第 5 章起）、惰性求值的代价（第 2 章）。

\section*{参考文献}
\addcontentsline{toc}{section}{参考文献}\label{ux53c2ux8003ux6587ux732e}

本章属于概念综述，没有依赖某一本书的推导。关键概念的出处与延伸资源如下：

\begin{itemize}
\tightlist
\item
  \textbf{Cook, William R.} “Object-Oriented Programming Versus Abstract Data Types.” 1990.
  讨论对象与抽象数据类型取舍的论文；本章 1.3.4 节的分析框架来自此文。
\item
  \textbf{Wadler, Philip.} “The Expression Problem.” Java Generics 邮件列表，1998.
  “表达问题”这一名称的出处；正文的两轴表可与该邮件中的表述对照。
\item
  \textbf{Awodey, Steve.} \emph{Category Theory}, 2nd ed. Oxford University Press, 2006. §1.3。
  本章未展开范畴论，但 1.3 节使用的对象与态射语言与该节定义一致。
\item
  \textbf{PureScript 文档}。“\href{https://github.com/purescript/documentation/blob/master/Related-Projects.md}{Related Projects}”。
  1.5 节对派生语言与同源语言的对比据此文档。
\item
  \textbf{Haskell 官网}。\url{https://www.haskell.org/}。下载、教程与社区入口，1.1 节的 GHC 与 GHCi
  都可以从这里获取。
\item
  \textbf{GHC 用户指南}。\url{https://downloads.haskell.org/ghc/latest/docs/users_guide/}，其中
  GHCi 一章给出 \passthrough{\lstinline!:t!}、\passthrough{\lstinline!:l!} 等命令的完整清单与 \passthrough{\lstinline!:\{ ... :\}!} 的规则。
\item
  \textbf{Hoogle}。\url{https://hoogle.haskell.org/}。按名字或类型检索 \passthrough{\lstinline!base!} 等包；查找 \passthrough{\lstinline!(.)!}、
  \passthrough{\lstinline!curry!}、\passthrough{\lstinline!uncurry!} 的签名用它最快。
\item
  \textbf{Python 官方文档}。\url{https://docs.python.org/3/}。1.1 与 1.4 节的 Python 对照代码与其
  语义一致，REPL 行为见官方教程。
\item
  \textbf{SICP}。Abelson 与 Sussman, \emph{Structure and Interpretation of Computer Programs},
  2nd ed. MIT Press, 1996；\href{https://mitp-content-server.mit.edu/books/content/sectbyfn/books_pres_0/6515/sicp.zip/index.html}{在线版}。
  1.5.1 节 Lisp 路线（代码即数据、宏）的系统叙述。
\end{itemize}
